\documentclass[conference]{IEEEtran}
\IEEEoverridecommandlockouts
\usepackage{cite}
\usepackage{amsmath,amssymb,amsfonts}
\usepackage{graphicx}
\usepackage{textcomp}
\usepackage{xcolor}

\usepackage{subcaption}

\usepackage[draft,hidelinks,bookmarks=false,bookmarksdepth=0]{hyperref}
\usepackage{xurl}

\usepackage{listings}

\def\BibTeX{{\rm B\kern-.05em{\sc i\kern-.025em b}\kern-.08em
    T\kern-.1667em\lower.7ex\hbox{E}\kern-.125emX}}
\begin{document}

\title{Carbon Footprint of Relational and Hybrid Object-Relational Storage for Image-Based Time-Series Workloads}

\author{
[Hidden for double-blind review]
}


\maketitle

\begin{abstract}
Smart building and facility-monitoring systems increasingly persist camera frames together with timestamped metadata, making storage-engine choice relevant not only for performance but also for operational carbon cost. This paper compares two architectures for image-based time-series workloads: PostgreSQL with TimescaleDB storing images inline as \texttt{BYTEA} (PG), and PostgreSQL with TimescaleDB externalising binaries to MinIO object storage while retaining metadata in the database (PM). Using isolated Docker deployments on an Ubuntu~24.04.3 host with an 11th Gen Intel Core i7-1165G7 processor, 15.3~GB RAM, and CodeCarbon machine-mode tracking, we benchmark seven quality-90 JPEG resolutions from 360p to 5K. The results show a resolution-dependent crossover. PG is faster for insertion from 360p through 1440p, while PM becomes faster at 4K and 5K. Latest-payload retrieval is nearly tied at the smallest profiles and PM is faster from 720p onward, reaching a $6.5\times$ latency advantage at 5K. Measured whole-phase totals show PM consuming 0.012286~kWh and emitting 8305~mg CO$_2$eq versus PG's 0.014912~kWh and 10080~mg CO$_2$eq, a 17.6\% reduction. Per-profile estimates show PM emits more at 360p--1080p but less from 1440p upward. These findings suggest inline \texttt{BYTEA} remains attractive for lower-resolution ingest-heavy workloads, whereas externalised object storage becomes preferable once payload size and retrieval sensitivity dominate.
\end{abstract}

\begin{IEEEkeywords}
time-series databases, IoT building monitoring, PostgreSQL, MinIO, object storage, carbon footprint, green computing
\end{IEEEkeywords}

\section{Introduction}
\label{sec:introduction}

Smart building and facility-monitoring deployments increasingly combine conventional telemetry with image streams from embedded cameras, doorbells, inspection devices, and other edge systems~\cite{ieee_ce_definition,intuz_iot_databases}. These applications often need to retain binary media together with timestamps, device identifiers, and other metadata so that operators can inspect the latest frame, replay recent events, or query activity over time~\cite{ceec_tsdb_review,jensen_tsms_survey}. In such systems, the data tier runs continuously and its energy use contributes directly to the operational carbon footprint of the building or edge infrastructure~\cite{nerc_tsdb_survey}.

Many organisations prefer to build on PostgreSQL~\cite{postgresql_official} and TimescaleDB~\cite{timescaledb_official} rather than introduce a specialised storage stack. When binary media is small, inline storage inside relational rows can be operationally convenient. When binary media becomes large, however, a common alternative is to keep only queryable metadata in PostgreSQL and store the bulk payload in an S3-compatible object store such as MinIO~\cite{minio_official} or Amazon S3~\cite{aws_s3_official}. These two options imply different write amplification, memory behaviour, and retrieval paths, and therefore potentially different energy profiles.

Existing benchmarks in the time-series literature focus primarily on lightweight numeric events~\cite{tsbs,cratedb_ts_benchmark,diva_timescaledb_mongodb_eval}. They rarely evaluate multi-megabyte image payloads, and even more rarely report energy or CO$_2$-equivalent emissions. This leaves a practical gap for image-heavy IoT deployments that must decide whether the simplicity of inline database storage outweighs the performance and sustainability benefits of separating metadata from binary objects.

This paper addresses the following research question: \textit{What are the performance and carbon-footprint trade-offs of inline PostgreSQL \texttt{BYTEA} storage (PG) versus PostgreSQL+MinIO externalised object storage (PM) for image-based time-series workloads across increasing image resolutions?}

To answer this question, we benchmark both architectures across seven JPEG profiles from 360p to 5K in isolated Docker containers, measure whole-phase energy and emissions with CodeCarbon~\cite{codecarbon}, and analyse insertion throughput, latest-payload point-read latency, storage amplification, and per-profile carbon estimates.

\section{Related Work}
\label{sec:related_work}

Time-series databases are designed to sustain append-oriented ingestion by combining time-based partitioning, compression, and storage layouts that reduce the write amplification of traditional B-Tree update paths~\cite{dewaal_timeseries_survey,jensen_tsms_survey}. TimescaleDB~\cite{timescaledb_official} extends PostgreSQL~\cite{postgresql_official} with hypertables while preserving SQL semantics and the operational ecosystem many organisations already use.

Most comparative benchmarks in this space target lightweight scalar data rather than binary media. TSBS~\cite{tsbs} uses small DevOps and telemetry records, CrateDB's benchmark~\cite{cratedb_ts_benchmark} focuses on numeric IoT metrics, and the evaluation by Mohamud and Brostr\"om~\cite{diva_timescaledb_mongodb_eval} studies numerical sensor workloads. These benchmarks do not characterise the behaviour of multi-megabyte image payloads, where storage-engine internals become far more consequential.

That gap matters because large binaries stress systems differently. PostgreSQL's TOAST~\cite{postgresql_toast_docs} mechanism relocates oversized values out of the main row, but the payload still remains within the database's transactional and retrieval machinery. Object stores such as MinIO~\cite{minio_official}, by contrast, are designed to stream large immutable objects directly while leaving relational systems responsible only for lightweight metadata and indexing. The architectural trade-off is therefore not simply ``one system versus two systems,'' but ``one transactional path for everything'' versus ``separating metadata management from binary transport.''

The carbon footprint of software systems has also attracted increasing attention. Lannelongue et al.~\cite{lannelongue_green_algorithms} formalised the Green Algorithms approach, and CodeCarbon~\cite{codecarbon} provides a practical implementation that estimates energy use and associated emissions from machine-level measurements and regional grid intensity. Within database systems, however, relatively little work applies this style of accounting to binary time-series storage workloads. Our study addresses that gap for relational versus hybrid object-relational image storage.

\section{Methodology}
\label{sec:methodology}

\subsection{Hardware and Software Environment}
All experiments run on a Lenovo ThinkPad L13 laptop running Ubuntu~24.04.3~LTS with Linux kernel 6.8.0-90-generic. The machine provides an 11th Gen Intel Core i7-1165G7 processor (4 cores / 8 threads, 2.80~GHz base, 4.70~GHz max), 15.3~GB RAM, and a 476.9~GB Samsung MZVL4512HBLU-00BL7 NVMe SSD. CodeCarbon detects the host in Indonesia (Jakarta region) and operates in \texttt{tracking\_mode=machine}, using host-level CPU and RAM measurements for the benchmark phases.

\begin{table}[htbp]
\centering
\renewcommand{\arraystretch}{1.1}
\caption{Software Stack}
\label{tab:software_versions}
\begin{tabular}{ll}
\hline
\textbf{Component} & \textbf{Version} \\
\hline
Docker          & 29.4.0 \\
Docker Compose  & v5.1.1 \\
PostgreSQL      & 15 + TimescaleDB (\texttt{latest-pg15}) \\
MinIO           & \texttt{minio/minio:latest} \\
Python          & 3.12.3 \\
psycopg2        & 2.9.11 \\
minio           & 7.2.20 \\
Pillow          & 12.1.0 \\
CodeCarbon      & 3.2.5 \\
\hline
\end{tabular}
\end{table}

\subsection{Environment and Workloads}
Both architectures are deployed in Docker containers on the same host. PostgreSQL exposes port~55432 and is launched with \texttt{max\_connections=100}. MinIO exposes ports~59000 and~59001 and uses the bucket \texttt{sensor-media}. The benchmark harness generates deterministic JPEG payloads from a single source image using a collage transform implemented in Python with Pillow. All profiles use JPEG quality~90.

Seven image profiles are evaluated: 360p ($640{\times}360$, $\sim$0.09~MB), 480p ($854{\times}480$, $\sim$0.16~MB), 720p ($1280{\times}720$, $\sim$0.36~MB), 1080p ($1920{\times}1080$, $\sim$0.80~MB), 1440p ($2560{\times}1440$, $\sim$1.40~MB), 4K~UHD ($3840{\times}2160$, $\sim$2.99~MB), and 5K ($5120{\times}2880$, $\sim$5.09~MB). For each profile, the benchmark performs a 100-row warm-up, followed by 5 measured insertion runs of 2,000 rows each in batches of 50. Latest-payload retrieval is measured as a point-read query with 2 warm-up runs and 5 measured runs. Driver overhead uses 10 warm-up and 10 measured roundtrips.

\subsection{Schema and Operations}
PG stores each frame as a TimescaleDB hypertable row with a \texttt{BYTEA} payload column, allowing PostgreSQL's TOAST~\cite{postgresql_toast_docs} mechanism to manage oversized values transparently. PM stores the metadata in PostgreSQL but replaces the inline payload with a \texttt{minio\_object\_key}; the binary image itself is uploaded to MinIO. Both tables use a compound index on \texttt{(device\_id, ts DESC)} to support latest-payload access.

Each measured insert run recreates the target table, and PM also recreates the MinIO bucket. For retrieval, PG returns the inline payload through PostgreSQL, whereas PM first retrieves the latest object key from PostgreSQL and then issues an HTTP GET to MinIO. This design isolates the core storage behaviour of each architecture while preserving realistic end-to-end operations.

\subsection{Carbon Footprint Measurement}
The benchmark runner wraps the PostgreSQL-only and PostgreSQL+MinIO phases separately with CodeCarbon. Each tracked phase therefore includes driver overhead, warm-up work, insert benchmarks, point-read benchmarks, and setup/teardown costs inside one measured duration. CodeCarbon reports total phase duration, energy consumed, CPU energy, RAM energy, and CO$_2$-equivalent emissions.

To obtain per-profile carbon estimates, we apply a constant power-rate model to the measured phase totals. For a phase with total energy $E$, emissions $C$, and duration $T$, we derive average rates $E/T$ and $C/T$ and multiply them by each profile's measured insert and point-read durations. These per-profile values are estimates rather than direct measurements, while the phase-level totals remain directly measured by CodeCarbon.

\begin{table*}[htbp]
\centering
\renewcommand{\arraystretch}{1.1}
\caption{Insert Throughput, Point-Read Latency, and Storage Amplification (2,000 rows)}
\label{tab:performance_summary}
\resizebox{\textwidth}{!}{%
\begin{tabular}{l rr rr rr}
\hline
\textbf{Res.} & \multicolumn{2}{c}{\textbf{Insert (rows/s)}} & \multicolumn{2}{c}{\textbf{Point Read (ms)}} & \multicolumn{2}{c}{\textbf{Storage Amp.}} \\
 & PG & PM & PG & PM & PG & PM \\
\hline
360p  & \textbf{424.23} & 189.52 & 1.777 & \textbf{1.338} & 1.0454 & \textbf{1.0029} \\
480p  & \textbf{407.13} & 74.71 & \textbf{1.253} & 1.329 & 1.0435 & \textbf{1.0017} \\
720p  & \textbf{164.05} & 61.37 & 2.615 & \textbf{1.524} & 1.0394 & \textbf{1.0007} \\
1080p & \textbf{71.95} & 51.16 & 4.932 & \textbf{1.631} & 1.0388 & \textbf{1.0003} \\
1440p & \textbf{39.60} & 36.91 & 7.358 & \textbf{2.195} & 1.0390 & \textbf{1.0002} \\
4K    & 16.72 & \textbf{24.28} & 17.563 & \textbf{3.082} & 1.0378 & \textbf{1.0001} \\
5K    & 8.55 & \textbf{10.85} & 27.208 & \textbf{4.167} & 1.0376 & \textbf{1.0001} \\
\hline
\multicolumn{7}{l}{\footnotesize PG = PostgreSQL+TimescaleDB with inline \texttt{BYTEA}. PM = PostgreSQL+TimescaleDB+MinIO.} \\
\end{tabular}
}
\end{table*}

\begin{table*}[htbp]
\centering
\renewcommand{\arraystretch}{1.05}
\caption{Per-Profile Carbon Footprint Estimates from the Constant Power-Rate Model}
\label{tab:carbon_summary}
\resizebox{\textwidth}{!}{%
\begin{tabular}{l rr rr rr rr rr}
\hline
\textbf{Res.}
  & \multicolumn{2}{c}{\textbf{Duration (s)}}
  & \multicolumn{2}{c}{\textbf{Energy (\(\mu\)Wh)}}
  & \multicolumn{2}{c}{\textbf{CPU (\(\mu\)Wh)}}
  & \multicolumn{2}{c}{\textbf{RAM (\(\mu\)Wh)}}
  & \multicolumn{2}{c}{\textbf{Emissions (mg CO$_2$eq)}} \\
  & PG & PM & PG & PM & PG & PM & PG & PM & PG & PM \\
\hline
360p  & 23.6 & 53.6 & 150.2 & 289.1 & 84.6 & 140.3 & 65.6 & 148.8 & 101.5 & 195.4 \\
480p  & 24.6 & 137.8 & 156.3 & 743.6 & 88.0 & 361.0 & 68.3 & 382.7 & 105.7 & 502.6 \\
720p  & 61.1 & 164.9 & 388.5 & 889.9 & 218.8 & 431.9 & 169.7 & 457.9 & 262.6 & 601.5 \\
1080p & 139.3 & 196.3 & 886.2 & 1059.5 & 499.1 & 514.3 & 387.0 & 545.2 & 599.0 & 716.1 \\
1440p & 252.7 & 272.8 & 1607.3 & 1472.5 & 905.3 & 714.8 & 702.0 & 757.7 & 1086.4 & 995.3 \\
4K    & 600.2 & 414.4 & 3817.0 & 2236.7 & 2149.9 & 1085.7 & 1667.1 & 1151.0 & 2580.0 & 1511.8 \\
5K    & 1171.4 & 924.5 & 7449.5 & 4990.3 & 4195.8 & 2422.3 & 3253.7 & 2568.0 & 5035.4 & 3373.1 \\
\hline
\end{tabular}
}
\vspace{2pt}
\begin{minipage}{\textwidth}
\footnotesize Measured whole-phase totals including driver overhead, warm-up, and setup: PG = 0.014912~kWh / 10079.6~mg CO$_2$eq; PM = 0.012286~kWh / 8304.8~mg CO$_2$eq.
\end{minipage}
\end{table*}

\section{Results}
\label{sec:results}

\begin{figure}[t]
\centering
\includegraphics[width=\columnwidth]{figures/boxplot_insert_throughput.pdf}
\caption{Insert throughput across 360p--5K. PG leads through 1440p, while PM overtakes at 4K and remains ahead at 5K.}
\label{fig:insert_throughput}
\end{figure}

\begin{figure}[t]
\centering
\includegraphics[width=\columnwidth]{figures/boxplot_retrieval_latency.pdf}
\caption{Latest-payload binary retrieval latency across 360p--5K. PG is slightly faster only at 480p; PM is faster for the remaining profiles and the gap widens sharply at 4K and 5K.}
\label{fig:retrieval_latency}
\end{figure}

\begin{figure}[t]
\centering
\includegraphics[width=\columnwidth]{figures/carbon_footprint.pdf}
\caption{Estimated per-profile CO$_2$eq emissions (mg) split into insert and retrieval contributions. PM emits more at small profiles but less from 1440p upward.}
\label{fig:carbon_footprint}
\end{figure}

\begin{figure*}[t]
\centering
\includegraphics[width=.9\textwidth]{figures/carbon_breakdown.pdf}
\caption{Per-profile energy breakdown (CPU + RAM, stacked bars) with benchmark duration shown as dashed lines. The crossover in PM's carbon advantage follows the crossover in end-to-end benchmark time.}
\label{fig:carbon_breakdown}
\end{figure*}

\subsection{Supporting Performance Observations}
\label{subsec:performance}
Table~\ref{tab:performance_summary} and Figures~\ref{fig:insert_throughput}--\ref{fig:retrieval_latency} show that insertion throughput exhibits a clear write-side crossover. PG is faster from 360p through 1440p, with especially large margins at 480p (407.13 vs.\ 74.71~rows/s, $5.45\times$) and 720p (164.05 vs.\ 61.37~rows/s, $2.67\times$). The gap narrows steadily as payload size increases and reverses at 4K, where PM reaches 24.28 vs.\ 16.72~rows/s ($1.45\times$ faster), and remains in PM's favour at 5K (10.85 vs.\ 8.55~rows/s, $1.27\times$ faster). On this host, inline \texttt{BYTEA} therefore remains a strong ingest option up to 1440p, but PM becomes more effective once payload sizes reach the largest profiles.

Latest-payload point-read latency follows a different pattern. PG is slightly faster only at 480p (1.253 vs.\ 1.329~ms), while PM is already faster at 360p (1.338 vs.\ 1.777~ms) and becomes decisively faster from 720p onward. At 1080p, PM reduces latency from 4.932 to 1.631~ms ($3.02\times$ faster); at 4K, from 17.563 to 3.082~ms ($5.70\times$ faster); and at 5K, from 27.208 to 4.167~ms ($6.53\times$ faster). The read-side crossover therefore occurs much earlier than the write-side crossover.

Storage efficiency is consistently in PM's favour. PG remains near $1.04\times$ storage amplification for all profiles, whereas PM stays essentially at raw-payload size, between $1.0029\times$ and $1.0001\times$. At 5K and 2,000 inserted rows, PG occupies 10555.0~MB versus 10172.9~MB for PM, a 382.1~MB gap despite similar logical content.

Driver overhead provides useful context. Averaged across the seven profiles, PostgreSQL's driver roundtrip measures 0.075~ms, while MinIO's \texttt{list\_buckets} roundtrip measures 0.469~ms. The MinIO path therefore carries a higher fixed per-request overhead, but this cost is quickly amortised once large image transfers dominate the workload.

\subsection{Carbon Footprint (RQ)}
\label{subsec:carbon_footprint}
The whole-phase CodeCarbon measurements show that PM is more carbon-efficient overall on this seven-profile benchmark. The PostgreSQL-only phase ran for 2344.8~s, consumed 0.014912~kWh, and emitted 10079.6~mg CO$_2$eq. The PostgreSQL+MinIO phase ran for 2276.1~s, consumed 0.012286~kWh, and emitted 8304.8~mg CO$_2$eq. This corresponds to a \textbf{17.6\% reduction} in both energy and emissions for PM at the phase level.

Table~\ref{tab:carbon_summary} and Figures~\ref{fig:carbon_footprint}--\ref{fig:carbon_breakdown} refine that picture by estimating the profile-level contribution of insert and retrieval work. These estimates show that PM is \emph{not} uniformly greener. At 360p, PM emits 195.4~mg CO$_2$eq versus 101.5~mg for PG (+92.5\%). At 480p, the PM penalty is much larger (502.6 vs.\ 105.7~mg, +375.7\%), reflecting its poor small-payload insert throughput. PM remains worse at 720p and 1080p, though the gap narrows substantially by 1080p (716.1 vs.\ 599.0~mg, +19.6\%).

The carbon crossover occurs at 1440p. There, PM emits 995.3~mg versus 1086.4~mg for PG, an 8.4\% advantage. At 4K the PM advantage becomes pronounced, dropping from 2580.0 to 1511.8~mg (-41.4\%), and at 5K the reduction remains substantial at 3373.1 versus 5035.4~mg (-33.0\%). This behaviour mirrors the throughput crossover observed in Section~\ref{subsec:performance}: once payload sizes become large enough for PostgreSQL's inline binary handling to dominate execution time, PM's separate object-storage path becomes both faster and lower-carbon.

The difference between the measured whole-phase totals and the summed profile estimates is expected. The per-profile table only allocates insert and retrieval time, whereas the tracked phases also include warm-up activity, driver-overhead measurements, and setup work. The directly measured phase totals therefore remain the correct basis for the overall RQ answer, while the profile estimates explain where that overall difference originates.

\section{Discussion}
\label{sec:discussion}

\textbf{Resolution-dependent crossover.} The updated benchmark does not support a universal ``always use MinIO'' or ``always use inline \texttt{BYTEA}'' conclusion. PG remains the better write engine through 1440p on this machine, even though PM already provides superior retrieval latency at most profiles. From 4K onward, PM becomes the better choice for both insertion throughput and point-read latency, and from 1440p onward it also becomes the lower-carbon option.

\textbf{Why PostgreSQL leads at smaller writes.} At lower resolutions, PG benefits from batched inserts that avoid an extra object-store request per payload. The images are small enough that PostgreSQL's inline binary handling and TOAST overhead remain tolerable, while PM still pays the fixed cost of a MinIO PUT plus a relational metadata write. That is why PM performs particularly poorly at 480p despite its much better storage amplification.

\textbf{Why PostgreSQL+MinIO wins at large payloads.} As payload size grows, the cost of keeping the binary within the relational engine rises sharply. PostgreSQL must still log, manage, and retrieve large values through the database path, whereas PM lets PostgreSQL manage only metadata while MinIO streams the object directly. This architectural separation turns into both a performance advantage and a carbon advantage once the binary payload is large enough, which in our measurements happens around 1440p for emissions and 4K for insertion throughput.

\textbf{Implications for system design.} These results suggest a workload-sensitive deployment rule. If the system is dominated by lower-resolution ingestion and operational simplicity is a priority, inline \texttt{BYTEA} storage remains defensible. If the system needs low-latency latest-frame access or expects 4K--5K payloads, externalising binaries to object storage becomes the more robust choice. PM also offers consistently better storage efficiency, which compounds as retention increases.

\textbf{Threats to validity.} Results are from a single laptop-class host and a single container topology. Different hardware, distributed object-storage deployments, or more compressible media could shift the exact crossover points. The per-profile carbon table is estimated from whole-phase measurements using a constant power-rate assumption, so it should be interpreted as a benchmark-scoped explanation rather than a hardware-instrumented per-operation trace. Even so, both architectures were benchmarked under identical conditions, making the comparative trends meaningful.

\section{Conclusions}
\label{sec:conclusions}

This study evaluated the performance and carbon footprint of two storage architectures for image-based time-series workloads: PostgreSQL with inline \texttt{BYTEA} storage and PostgreSQL with externalised MinIO object storage.

\textbf{RQ answer:} The better architecture depends on payload size and workload emphasis, but the complete seven-profile benchmark on this machine favours PostgreSQL+MinIO overall. Measured whole-phase totals show PM emitting 8304.8~mg CO$_2$eq versus PostgreSQL's 10079.6~mg, a \textbf{17.6\% reduction}. PG remains faster for insertion from 360p through 1440p, while PM becomes faster at 4K and 5K. For latest-payload retrieval, PM is faster at every profile except a slight PostgreSQL advantage at 480p, and the gap reaches $6.5\times$ at 5K.

The key practical finding is therefore a pair of crossover points rather than a universal winner. On this host, inline PostgreSQL is attractive for smaller ingest-heavy workloads, while PostgreSQL+MinIO becomes the more compelling design as binary payload size increases and retrieval or carbon efficiency become dominant concerns. Future work should evaluate the same architectures on multi-node deployments and under higher client concurrency to determine how stable these crossover points remain outside a single-host setup.

\section*{Acknowledgment}
[Hidden for double-blind review]


\bibliographystyle{IEEEtran}
\bibliography{references}

\end{document}
